\section{Mistral Large 4のオープンウェイト公開が話題に}
Mistralの大規模モデル「Large 4（Le Chonk）」の重み公開が、観測期間中の複数のAIニュース投稿で紹介された。投稿では総パラメータ約1.05兆、活性化約520億、最大100万トークンの文脈長、マルチモーダル対応を掲げる。エージェントコーディングやサイバー防御向けの性能向上も主張されている。ただし、公式発表は観測窓外とされ、ライセンスの商用ホスト条件やベンチマークの優位性は本号で独立確認していない。今回は公開への反響を記録する。
\dailyxsource{Maha (@mahadaajani)、観測窓内の紹介}{https://x.com/mahadaajani/status/2116289797006500266}

\section{Gemini 4 Carbonのサイバー防御性能を議論}
GoogleのGemini 4 Carbonについて、サイバーセキュリティ向けの限定提供と評価スコアがXで再び取り上げられた。投稿ではArtificial Analysis Intelligence Indexの59点、SEC-bench Proの88.2\%などが紹介され、Fairwind Program経由でのアクセスが説明される。発表自体は10月28日頃とされ、今回の観測窓より前である。GrokはGoogleの一次資料を直接確認しておらず、順位や性能を確定的な比較結果とは扱わない。一般開発者への提供範囲も今後の確認事項となる。
\dailyxsource{田嶋正吾 (@st\_dev0)、窓内の紹介}{https://x.com/st_dev0/status/2116289893064462387}

\section{Project Suncatcher、宇宙TPUから応答を送信}
GoogleはProject Suncatcherのプロトタイプ衛星でTPUを動かし、プロンプトへの応答として俳句を地球に返したと公式Xで発表した。宇宙環境でのAI推論を実証する試みとして注目され、2027年には相互接続する衛星の計画も示されている。今回の記録では、実験の詳細な電力条件、使用モデルや長時間運用の性能までは確定していない。単発の応答成功と、継続的な宇宙データセンター運用の実現性は区別して評価する必要がある。
\dailyxsource{Google (@Google)}{https://x.com/Google/status/2116244898269610456}

\section{Genie 4、研究プレビューの提供状況を更新}
Googleはインタラクティブなワールドモデル「Genie 4」の研究プレビューを、信頼できるテスターに展開していると公式Xで案内した。ネイティブ音声対応のデモも紹介し、今後はGemini APIパートナーやGoogle AI Ultraへの提供を見込むという。モデル自体の主要発表は観測窓より前だが、今回の公式投稿は提供状況に関する更新に当たる。一般公開の時期やAPI仕様、エージェント学習への応用は今後の確認事項である。
\dailyxsource{Google (@Google)}{https://x.com/Google/status/2116237215026032872}

\section{OpenAIのモデル監視と行動仕様が議論に}
OpenAIのフロンティアモデル監視システム拡充とModel Specの五原則への再編が、ニュースまとめ投稿などで紹介された。エージェントのサンドボックス脱出などの兆候を検知し、過去のインシデントの再現実験では早期警告が可能だったとするが、これはOpenAI側の報告であり独立検証はない。Grokが挙げた公式X投稿はAIの未来を扱う動画で、監視システムの詳細を直接立証する投稿とは区別が必要だ。監視対象や仕様差分は追加資料待ちとなる。
\dailyxsource{OpenAI (@OpenAI)、公式動画投稿}{https://x.com/OpenAI/status/2116253388883345623}

\section{カナダ政府によるCohere出資の報道}
カナダ政府がCohereに最大16億カナダドルを出資し、国内の本社維持と計算資源への投資を条件とするとの情報がXで議論された。Mark Carney首相の投稿を引用した賛否も広がったという。ただし、首相の元投稿は10月29日で観測窓外に当たり、政府発表の一次資料も直接確認できていない。金額と条件は投稿上の主張として扱い、正式な契約内容と区別する。
\dailyxsource{Mark Carney (@MarkJCarney)、元投稿は観測窓外}{https://x.com/MarkJCarney/status/2115946856563741084}

\section{Anthropic、AI for Scienceの支援拡大が話題に}
Anthropicの「AI for Science」が20の科学機関との協力に拡大し、Claudeの利用クレジットや研究協力を提供するとする情報がXで共有された。Allen Instituteなどとの研究活動も紹介され、AIモデルを科学研究の実務へ組み込む動きとして注目された。Grokは公式ニュースの存在を記録しているが、発表日時や内容を直接照合していない。実際の共同研究成果や支援条件は今後の確認事項であり、今回は観測窓内の紹介投稿を中心に整理した。
\dailyxsource{Maha (@mahadaajani)}{https://x.com/mahadaajani/status/2116158934050582562}

\section{ChatGPT Proなどの利用上限変更を巡る投稿}
ChatGPTのPro 100／200／500およびBusiness Premiumで、WorkやCodexの利用上限が倍増したとの投稿が共有された。従来のPro 200向け削減計画が撤回されたという説明も含まれる。ただし、Grokの収集ではOpenAI公式の告知やヘルプセンターとの照合がなく、料金プラン別の正確な条件は未確認である。利用上限は契約や時期により変わり得るため、現段階では投稿者による変更報告として扱い、確定したサービス仕様とはみなさない。
\dailyxsource{内涵段子 (@hacken2006)、未確認の変更報告}{https://x.com/hacken2006/status/2116290647975981329}

\section{Canvaのデザイン向けモデルBrushが話題に}
Canvaがデザインの内容を理解して編集を支援するビジョン言語モデル「Brush」を導入し、Canva AIで利用可能にしたとの情報がXのニュースまとめで紹介された。生成AIを単に画像作成へ用いるだけでなく、既存デザインの構造を理解した編集へ広げる方向性を示す話題である。ただし、GrokはCanvaの公式発表を直接確認できておらず、モデルの技術仕様、提供範囲、編集機能の詳細は不明である。今回は紹介投稿の内容にとどめる。
\dailyxsource{AlexAImaginator (@TraffAlex)}{https://x.com/TraffAlex/status/2116168604924227763}

\section{DeepgramによるDaily買収の情報が共有}
音声認識を手掛けるDeepgramが、音声エージェント向けフレームワーク「Pipecat」を持つDailyを買収したとの報告が、観測窓内のAIニュース投稿で共有された。音声認識から対話エージェントの開発基盤までを一体化する動きとして関心を集める。ただし、Grokは公式プレスリリースや取引条件を直接確認していない。買収の詳細、Pipecatの今後の開発・公開体制については、正式な発表を待つ必要がある。
\dailyxsource{Kana (@kana93\_ai)}{https://x.com/kana93_ai/status/2116290372808896938}

\section{Hy-Nomas、Genie 4との比較動画を公開}
研究者のHeng Wang氏が、インタラクティブなオープンワールドモデル「Hy-Nomas」を紹介し、Genie 4との比較動画をXに投稿した。投稿ではオープンソース化やコンシューマ向けGPUでの動作を主張している。ただし、GrokはGitHubの正確なリポジトリURLや技術レポート、モデル重みを検証しておらず、性能と公開範囲は未確認である。比較動画は研究上の方向性を示す材料となるが、条件をそろえた性能評価とは区別して読む必要がある。
\dailyxsource{Heng Wang (@hengwang00)、性能・公開状況は未検証}{https://x.com/hengwang00/status/2116267208628437501}

\section{Reddit、AI検索による流入減少を巡る報道}
RedditのCEOが、Google AI Modeの影響で週100万件以上の検索経由リファラルが減少したと述べ、ログイン必須化を計画しているとの報道がX上で共有された。AI検索による要約が外部サイトへの流入を減らす可能性と、プラットフォーム側の対抗策が議論の焦点となっている。ただし、Grokが確認したのはニュースまとめと報道への言及であり、CEO本人の一次投稿は特定されていない。流入減少の原因や施策の詳細は未検証である。
\dailyxnote{Primary X投稿の直接URLは原文で未特定のため掲載しない。}
